\section{Stand der Forschung}
Die Personalisierung von Chatbots stellt gerade im Unternehmenskontext eine aktuelle Herausforderung dar. Ziel dieses Abschnittes ist es in diesem Zusammenhang aktuelle Techniken und Ansätze aus der Forschung vorzustellen, die sich ebenfalls mit dieser Fragestellung auseinandergesetzt haben.

\subsection{Personalisierte LLMs}
Ein oft behandeltes Thema bilden personalisierte LLMs. Dabei geht es darum LLMs explizit auf die Präferenzen eines Nutzers abzustimmen, sodass diese antworten generieren, die direkt auf eine Person zugeschnitten sind. In der Praxis stellt dies jedoch eine grundlegende Schwierigkeit für Unternehmen dar, da moderne LLMs meist in den Händen der Anbieter sind und sich nicht einfach so neu trainieren lassen. Aus diesem Grund bildet dieser Ansatz, sofern nicht eigene LLMs aufgebaut werden, einen unrealistischen Ansatz für diese Arbeit.

\subsection{Chat Historie}
Ein weiterer Ansatz der in vielen Arbeiten erwähnt wird, ist es die Chat Historie als Kontext in den Prompt zu integrieren. Das LLM könnte in diesem Zusammenhang selbst, die Präferenzen aus der Chat Historie extrahieren und anhand dieser eine Antwort generieren. In Anbetracht der Größe des Kontextfensters für LLMs, zeigen sich dabei jedoch grundlegende Herausforderungen, da die Anzahl an Tokens begrenzt ist. 

\subsection{RAG Ansätze}
Im Jahr 2025 stellten Zhao et al. in ihrem wissenschaftlichen Paper "Do LLMs Recognize Your Preferences? Evaluating Personalized Preference Following in LLMs" einen Benchmark (PrefEval) zur Evaluation der Präferenzverarbeitung in konversationsbasierten LLM-Systemen vor. In dieser Arbeit wurde untersucht, wie gut LLM Modelle Nutzerpräferenzen über längere Konversationen hinweg erkennen, speichern und in Antworten berücksichtigen können.

Die Ergebnisse des Benchmarks zeigen, dass moderne LLMs grundsätzlich in der Lage sind, Nutzerpräferenzen innerhalb von Konversationen zu berücksichtigen. Gleichzeitig zeigt sich jedoch, dass insbesondere die konsistente Verarbeitung von Präferenzen über längere und mehrere Konversationen hinweg weiterhin eine zentrale Herausforderung darstellt. 

Die Ergebnisse des Benchmarks zeigen, dass die Fähigkeit von LLMs zur Berücksichtigung von Nutzerpräferenzen mit zunehmender Gesprächslänge deutlich abnimmt. Während Präferenzen in kurzen Konversationen häufig noch korrekt verarbeitet werden können, sinkt die Genauigkeit bei längeren Dialogen erheblich. Insbesondere Präferenzen, die früh innerhalb einer Konversation genannt wurden, werden von den Modellen häufig nicht mehr zuverlässig berücksichtigt.

%Die Arbeit identifiziert dabei mehrere Problemfelder moderner LLMs im Kontext personalisierter Interaktionen. Dazu zählen unter anderem sogenannte Preference-Unaware Violations, bei denen Nutzerpräferenzen vollständig ignoriert werden, sowie Preference Hallucinations, bei denen Modelle fälschlicherweise Präferenzen ergänzen oder annehmen, die zuvor nicht geäußert wurden. Diese Ergebnisse verdeutlichen, dass große Kontextfenster allein nicht ausreichen, um eine konsistente Personalisierung sicherzustellen. Obwohl relevante Informationen technisch weiterhin im Kontext enthalten sein können, gelingt es den Modellen häufig nicht, diese Informationen zuverlässig abzurufen und in die Antwortgenerierung einzubeziehen.

Die Ergebnisse von PREFEVAL zeigen jedoch auch, dass insbesondere RAG-basierte Ansätze Halluzinationen reduzieren und die konsistente Berücksichtigung von Nutzerpräferenzen verbessern können. Durch den gezielten Abruf relevanter Präferenzinformationen aus externen Speicherstrukturen gelingt es den Modellen besser, frühere Nutzerangaben in späteren Konversationen erneut zu berücksichtigen. Dies ist insbesondere für personalisierte Chatbots im Unternehmenskontext relevant, da Nutzerpräferenzen häufig über längere Zeiträume hinweg konsistent verarbeitet werden müssen.

Externe Retrieval-Mechanismen bilden somit einen vielversprechenden Ansatz, um die Limitierungen kontextbasierter Präferenzverarbeitung moderner LLMs zu adressieren.